\section{Results}

\subsection{Token-ID transcription}

On development folds 4--6, \ImitatorEOne{} obtains 18.1, 19.8, and 16.9\% strict
accuracy, versus 2.2, 2.3, and 1.8\% for the initial autoregressive model.  Direct
token-ID prediction is therefore viable but not solved: \ImitatorEOne{} improves
strict sequence accuracy by roughly an order of magnitude over the initial
autoregressive decoder,
yet most complete sequences still contain at least one error.  The positional
intervention changes pairwise label order from near chance to approximately 0.99,
supporting the claim that the emitted sequence is visually ordered.

\begin{table}[t]
\centering\scriptsize
\caption{Prospectively registered signer-held-out confirmation.  The CIF column
is the frozen historical comparator, not a candidate \system{} architecture.}
\label{tab:confirmation}
\begin{tabular}{rrrrrrrr}
\toprule
Fold & CIF exact & \ImitatorEOne{} exact & \ImitatorEOne{} edit & Order & $\Delta_{\mathrm{perm}}$ & W/L/T & $p$ \\
\midrule
7 & \FoldSevenCIFExact & \FoldSevenExact & \FoldSevenEdit & \FoldSevenOrder &
    \FoldSevenPermutationDrop & \FoldSevenWLT & \FoldSevenPValue \\
8 & \FoldEightCIFExact & \FoldEightExact & \FoldEightEdit & \FoldEightOrder &
    \FoldEightPermutationDrop & \FoldEightWLT & \FoldEightPValue \\
\bottomrule
\end{tabular}
\end{table}

Under the registered conjunctive decision rule, H1 \HOneOutcome, H2
\HTwoOutcome, and H3 \HThreeOutcome.  Overall confirmation therefore
\OverallOutcome.  These statements are populated automatically from the frozen
evaluation artifacts and are not manually reinterpreted after either fold.

\begin{figure}[t]
\centering
\includegraphics[width=\textwidth]{confirmation_results.pdf}
\caption{Confirmatory evidence on unseen signers. Left: strict sequence accuracy
against the historical CIF comparator. Right: edit similarity collapses when
complete sign segments are permuted while targets remain fixed.}
\label{fig:confirmation-results}
\end{figure}

\begin{figure}[t]
\centering
\includegraphics[width=\textwidth]{keypoint_interpretability.pdf}
\caption{Qualitative keypoint interpretation for the target subtokens of an
unseen signer-7 ``A tierra / Perfume'' sequence. Left: normalized
Input$\times$Gradient attribution aggregated into anatomical regions. Right:
decrease in target log-probability when each region is zeroed. Both analyses
identify the left hand as the dominant evidence for the central subtokens; this
single example is descriptive rather than population-level evidence.}
\label{fig:keypoint-interpretability}
\end{figure}

Figure~\ref{fig:keypoint-interpretability} shows that the left hand contributes
28.9--41.6\% of normalized attribution and causes the largest ablation loss for
the central ``a'' and ``perf'' decisions (7.92 and 15.77 log-probability points).
Face keypoints also receive substantial local attribution, whereas their ablation
effect is smaller; this distinction indicates sensitivity without equivalent
causal necessity.  Pose matters mainly for ``perf,'' and the final ``ume'' subtoken
is comparatively insensitive to all four coarse removals.  The agreement between
hand saliency and hand ablation supports a manually grounded decision in this
example, but one selected sequence cannot establish a dataset-wide anatomy claim.

\subsection{Can \Gemma{} repair token errors?}

Table~\ref{tab:gemma-tolerance} answers this question for the current isolated-label
regime.  At zero corruption, \Gemma{} preserves exact content in 95\% of rows and
keeps chrF near 100.  At every tested nonzero rate, exact accuracy is already zero
before correction and remains zero afterward.  chrF decreases rather than improves.
Thus, no nonzero error tolerance is demonstrated: there is no tested nonzero
corruption level at which the language stage recovers exact content.  Because the
short sequences quantize every nonzero nominal rate to at least one substitution,
this experiment does not locate a continuous threshold between 0 and 0.10.

\begin{table}[t]
\centering\small
\caption{\Gemma{} 3n E2B correction audit on 120 isolated-label rows per
condition.  chrF is on a 0--100 scale.  On these short sequences, every nonzero
nominal rate produces at least one substitution.}
\label{tab:gemma-tolerance}
\begin{tabular}{rrrrrr}
\toprule
Substitution rate & Exact before & Exact after & chrF before & chrF after & $\Delta$chrF \\
\midrule
0.00 & 1.000 & 0.950 & 100.00 & 99.16 & -0.84 \\
0.10 & 0.000 & 0.000 & 36.67 & 34.58 & -2.10 \\
0.20 & 0.000 & 0.000 & 36.67 & 34.58 & -2.10 \\
0.30 & 0.000 & 0.000 & 36.67 & 34.58 & -2.10 \\
0.40 & 0.000 & 0.000 & 34.29 & 32.35 & -1.94 \\
0.50 & 0.000 & 0.000 & 23.30 & 21.85 & -1.45 \\
\bottomrule
\end{tabular}
\end{table}

This is not evidence that language models cannot correct sign transcription in
general.  A substitution such as ``photo'' $\rightarrow$ ``green'' in an isolated
label is compatible with no surrounding syntax or semantics that identifies the
original word.  The appropriate next experiment is therefore not another prompt
sweep on LSA64, but a pre-registered sentence corpus in which controlled token
errors leave genuine contextual redundancy.
